\section{Circuit-Level Diagnostic Run}
\label{app:diagnostic}

For transparency we report the circuit-level-only configuration, which disables the five program-level relation families and instead applies the four classic \emph{circuit-level} transformations---identity insertion, swap-based rewriting, phase-preserving rewriting, and circuit-equivalence rewriting---to the final circuit. We run it on a circuit-level Bugs4Q set (32 subjects) and an injected set (250 variants), with MorphQ matched on the same subjects, backends, shots, seeds, and timeouts. To reconcile denominators: the 32 are all Bugs4Q programs that build and run under our pinned Qiskit, of which the program-level audit (Sec.~\ref{sec:rq1}) admits the 14 that carry a candidate program relation; the 250 circuit-level injected variants and the 100 program-level injected mutants are distinct benchmarks from different operators. It shows \emph{no} advantage over MorphQ: $2/32$ on Bugs4Q and $17/250$ on the injected set, \emph{identical} for both methods. The per-transformation ablation attributes all detection to swap (\texttt{swap}: $17/250$; \texttt{identity}, \texttt{phase}, \texttt{equivalence}: $0/250$ each; any combination containing \texttt{swap}: $17/250$). Because a circuit-level transformation rewrites the final circuit a builder produced, it cannot expose a defect already encoded there---which is exactly the motivation for defining relations over program inputs. These numbers are diagnostic only and must not be quoted as program-level results.
